<<<<<<< HEAD
\documentclass[12pt,a4paper]{article}

\usepackage[T2A]{fontenc}
\usepackage[utf8]{inputenc}
\usepackage[english,russian]{babel}
\usepackage{amsmath,amssymb,amsthm}
\usepackage{graphicx}
\usepackage[hidelinks]{hyperref}
\usepackage[margin=2.5cm]{geometry}

\title{Название работы}
\author{К. И. Дик}
\date{\today}

\begin{document}

\maketitle

\section{Введение}

Начните писать здесь.

\begingroup
\renewcommand{\refname}{Список использованных источников}
\let\section\subsection
\addcontentsline{toc}{subsection}{\refname}
\begin{thebibliography}{9}

\bibitem{semenov-soprunov}
Семёнов, А.~Л. Решётка определимости (редуктов) для целых чисел с операцией
следования / А.~Л. Семёнов, С.~Ф. Сопрунов. --- Текст : электронный~//
Известия Российской академии наук. Серия математическая. --- 2021. ---
Т.~85, №~6. --- С.~245--258. --- DOI~10.4213/im9107. ---
URL: \url{https://www.mathnet.ru/rus/im9107}
(дата обращения: 19.08.2026).

\bibitem{libkin-fmt}
Libkin, L. Elements of Finite Model Theory / L.~Libkin. ---
Berlin ; Heidelberg : Springer, 2004. --- Текст : электронный. ---
DOI~10.1007/978-3-662-07003-1. ---
URL: \url{https://homepages.inf.ed.ac.uk/libkin/fmt/fmt.pdf}
(дата обращения: 19.08.2026).

\end{thebibliography}
\endgroup

\end{document}
=======
\documentclass[12pt,a4paper]{article}

\usepackage[T2A]{fontenc}
\usepackage[utf8]{inputenc}
\usepackage[english,russian]{babel}
\usepackage{amsmath,amssymb,amsthm}
\usepackage{graphicx}
\usepackage[hidelinks]{hyperref}
\usepackage[margin=2.5cm]{geometry}

% Unicode symbols retained from the original research notes.
\DeclareUnicodeCharacter{00B0}{\ensuremath{^{\circ}}}
\DeclareUnicodeCharacter{00B1}{\ensuremath{\pm}}
\DeclareUnicodeCharacter{00B2}{\ensuremath{^{2}}}
\DeclareUnicodeCharacter{00B7}{\ensuremath{\cdot}}
\DeclareUnicodeCharacter{03B4}{\ensuremath{\delta}}
\DeclareUnicodeCharacter{03B5}{\ensuremath{\varepsilon}}
\DeclareUnicodeCharacter{03C6}{\ensuremath{\varphi}}
\DeclareUnicodeCharacter{1D62}{\ensuremath{_{i}}}
\DeclareUnicodeCharacter{207A}{\ensuremath{^{+}}}
\DeclareUnicodeCharacter{207B}{\ensuremath{^{-}}}
\DeclareUnicodeCharacter{2081}{\ensuremath{_{1}}}
\DeclareUnicodeCharacter{2082}{\ensuremath{_{2}}}
\DeclareUnicodeCharacter{2098}{\ensuremath{_{m}}}
\DeclareUnicodeCharacter{209B}{\ensuremath{_{s}}}
\DeclareUnicodeCharacter{21D4}{\ensuremath{\Leftrightarrow}}
\DeclareUnicodeCharacter{2203}{\ensuremath{\exists}}
\DeclareUnicodeCharacter{2208}{\ensuremath{\in}}
\DeclareUnicodeCharacter{2261}{\ensuremath{\equiv}}
\DeclareUnicodeCharacter{2264}{\ensuremath{\le}}
\DeclareUnicodeCharacter{2265}{\ensuremath{\ge}}
\DeclareUnicodeCharacter{2C7C}{\ensuremath{_{j}}}

\providecommand{\tightlist}{%
  \setlength{\itemsep}{0pt}\setlength{\parskip}{0pt}}

\title{Название работы}
\author{К. И. Дик}
\date{\today}

\begin{document}

\maketitle

\section{Исследовательские заметки}

\subsection{Контекст проекта}

Долгосрочная тема исследования --- полная решётка определимости структуры
рациональных чисел со сложением и возможность автоматизации отдельных этапов
её построения.

Существующая работа задаёт контекст исследования. В ней собраны известные
результаты, проведённые эксперименты и исследованная часть решётки, однако не
предполагается, что в ней уже содержится метод получения полной решётки.

В работе рассматривались перебор формул и проверки с помощью SMT-решателя.
Отмеченные ограничения этих подходов: быстрый рост пространства формул,
сложность задач с кванторами, поиск функций как задача логики второго порядка,
а также наличие бесконечных семейств в решётке.

\subsection{Текущий объект исследования}

Первый этап исследования --- попытка построить решётку определимости
неориентированной квадратной решётки

\[
G=\langle\mathbb Z^2,E\rangle,
\]

где две точки смежны тогда и только тогда, когда расстояние между ними по
клеткам равно единице:

\[
E((x_1,y_1),(x_2,y_2))
\iff
|x_1-x_2|+|y_1-y_2|=1.
\]

Если \(R\) обозначает исходное ориентированное движение вправо или вверх, то

\[
E(p,q)\iff R(p,q)\lor R(q,p).
\]

Следовательно, \(E\) определимо через \(R\). Любое отношение, определимое через
\(E\), также определимо через \(R\). Поэтому решётка определимости
\(\langle\mathbb Z^2,E\rangle\) является подрешёткой исходной решётки --- главным
идеалом, расположенным ниже элемента \(E\) в порядке определимости.

\subsection{Соглашение о координатах}

Для точки со старыми координатами \((x,y)\) используются новые координаты

\[
u=x+y,
\qquad
v=x-y.
\]

Первая координата \(u\) является номером диагонали \(x+y=u\).

Вторая координата \(v\) задаёт положение относительно прямой \(x=y\). Знаковое
евклидово расстояние до этой прямой равно

\[
\frac{x-y}{\sqrt2}=\frac{v}{\sqrt2}.
\]

Обратное преобразование имеет вид

\[
x=\frac{u+v}{2},
\qquad
y=\frac{u-v}{2}.
\]

Если \((x,y)\in\mathbb Z^2\), координаты \((u,v)\) удовлетворяют условию

\[
u\equiv v\pmod2.
\]

Следовательно, образом целочисленной решётки является подрешётка индекса два:

\[
L=\{(u,v)\in\mathbb Z^2:u\equiv v\pmod2\}.
\]

В частности,

\[
(0,0)\mapsto(0,0),
\qquad
(1,0)\mapsto(1,1),
\qquad
(0,1)\mapsto(1,-1).
\]

\subsection{Сдвиги в новых координатах}

Сдвиг старых координат на вектор \((a,b)\),

\[
(x,y)\mapsto(x+a,y+b),
\]

в новых координатах принимает вид

\[
(u,v)\mapsto(u+a+b,v+a-b).
\]

Таким образом, старые базисные направления переходят в

\[
(1,0)\mapsto(1,1),
\qquad
(0,1)\mapsto(1,-1).
\]

Диагональное и антидиагональное направления становятся координатными осями:

\[
(k,k)\mapsto(2k,0),
\qquad
(k,-k)\mapsto(0,2k).
\]

\subsection{Автоморфизмы в новых координатах}

Перестановка старых координат

\[
(x,y)\mapsto(y,x)
\]

становится отражением

\[
(u,v)\mapsto(u,-v).
\]

Преобразование

\[
(x,y)\mapsto(-y,-x)
\]

становится отражением

\[
(u,v)\mapsto(-u,v).
\]

Поворот

\[
(x,y)\mapsto(y,-x)
\]

принимает вид

\[
(u,v)\mapsto(-v,u).
\]

Общий параллельный перенос

\[
(x,y)\mapsto(x+i,y+j)
\]

принимает вид

\[
(u,v)\mapsto(u+i+j,v+i-j).
\]

Все эти преобразования сохраняют условие одинаковой чётности координат:

\[
u\equiv v\pmod2.
\]

\subsection{Открытые вопросы}

\begin{itemize}
\tightlist
\item
  Как выглядят все классы с точностью до взаимной определимости?
\item
  Какие свойства конфигурации рациональных линейных подпространств определяют
  класс взаимной определимости?
\item
  Какие симметрии подрешётки
  \(L=\{(u,v)\in\mathbb Z^2:u\equiv v\pmod2\}\) продолжаются до автоморфизмов
  исходной бесконечной структуры?
\item
  Как меняется двумерное описание, если выбранные множители зависимы или
  порождают решётку ранга, отличного от двух?
\item
  Какие наблюдения для конечных факторов продолжаются до символически заданных
  автоморфизмов бесконечной структуры?
\end{itemize}

\section{\texorpdfstring{Критерий достижимости точки через \texttt{D(n,k)}}{Критерий достижимости точки через D(n,k)}}

Пусть \texttt{n} и \texttt{k} --- положительные целые числа и

\begin{verbatim}
d = gcd(n,k),
p = n/d,
q = k/d.
\end{verbatim}

Один шаг \texttt{D(n,k)} изменяет точку на один из векторов

\begin{verbatim}
(±n,±k) или (±k,±n),
\end{verbatim}

где знаки выбираются независимо.

Например, шаг с вектором \texttt{(n,k)} означает переход

\begin{verbatim}
(X,Y) -> (X+n,Y+k).
\end{verbatim}

Здесь \texttt{(n,k)} --- упорядоченная пара, то есть вектор перемещения. Это не
обозначение наибольшего общего делителя. Наибольший общий делитель везде
записывается как \texttt{gcd(n,k)} или как \texttt{d}.

Если сначала пройти путь с общим перемещением \texttt{U}, а затем путь с общим
перемещением \texttt{V}, результатом будет перемещение \texttt{U+V}. Противоположный
вектор для каждого разрешённого шага также разрешён. Поэтому любой уже
построенный путь можно пройти в обратную сторону и получить перемещение
\texttt{-U}.

\subsection{Теорема}

Точка \texttt{(x,y)} достижима из \texttt{(0,0)} через конечное число шагов \texttt{D(n,k)}
тогда и только тогда, когда выполняются следующие условия:

\begin{enumerate}
\def\labelenumi{\arabic{enumi}.}
\tightlist
\item
  \texttt{d} делит \texttt{x};
\item
  \texttt{d} делит \texttt{y};
\item
  если \texttt{p} и \texttt{q} оба нечётные, то \texttt{x/d} и \texttt{y/d} имеют одинаковую
  чётность.
\end{enumerate}

Если одно из чисел \texttt{p}, \texttt{q} чётное, третьего условия нет.

\subsection{Лемма Безу}

Существуют целые числа \texttt{a} и \texttt{b}, для которых

\begin{verbatim}
d = an + bk.
\end{verbatim}

\subsection{Необходимость условий}

В каждом шаге обе координаты изменяются на \texttt{±n} или \texttt{±k}. Оба числа
делятся на \texttt{d}. Поэтому после любого количества шагов обе координаты
также делятся на \texttt{d}. Следовательно,

\begin{verbatim}
d | x,
d | y.
\end{verbatim}

Теперь предположим, что \texttt{p} и \texttt{q} оба нечётные. После деления координат
каждого шага на \texttt{d} получается один из векторов

\begin{verbatim}
(±p,±q) или (±q,±p).
\end{verbatim}

Обе координаты такого вектора нечётные. Поэтому каждый шаг одновременно
меняет чётность обеих координат. В точке \texttt{(0,0)} координаты имеют
одинаковую чётность, следовательно, это остаётся верным во всех
достижимых точках:

\begin{verbatim}
x/d ≡ y/d (mod 2).
\end{verbatim}

Необходимость доказана.

\subsection{Достаточность условий}

Сначала покажем, что всегда достижимы векторы \texttt{(2d,0)} и \texttt{(0,2d)}.

Сделаем подряд два разрешённых шага: сначала \texttt{(n,k)}, затем \texttt{(n,-k)}.
Их изменения координат складываются:

\begin{verbatim}
(n,k) + (n,-k) = (2n,0).
\end{verbatim}

Поэтому \texttt{(2n,0)} --- перемещение некоторого пути из двух шагов. Аналогично,

\begin{verbatim}
(k,n) + (k,-n) = (2k,0).
\end{verbatim}

По лемме Безу существуют целые \texttt{a,b}, для которых \texttt{d\ =\ an\ +\ bk}.
Повторим путь с перемещением \texttt{(2n,0)} \texttt{\textbar{}a\textbar{}} раз и путь с перемещением
\texttt{(2k,0)} \texttt{\textbar{}b\textbar{}} раз. Если коэффициент отрицательный, соответствующий путь
проходится в обратную сторону. Суммарное перемещение равно

\begin{verbatim}
a(2n,0) + b(2k,0)
= (2an+2bk,0)
= (2d,0).
\end{verbatim}

Значит, \texttt{(2d,0)} достижима. Для второй координаты два разрешённых шага
дают

\begin{verbatim}
(n,k) + (-n,k) = (0,2k),
(k,n) + (-k,n) = (0,2n),
\end{verbatim}

а та же комбинация Безу даёт

\begin{verbatim}
a(0,2n) + b(0,2k) = (0,2d).
\end{verbatim}

Таким образом, достижимы \texttt{(2d,0)}, \texttt{(0,2d)}, их целые кратные и суммы.

\subsubsection{\texorpdfstring{Случай 1: одно из чисел \texttt{p}, \texttt{q} чётное}{Случай 1: одно из чисел p, q чётное}}

Так как \texttt{p} и \texttt{q} взаимно просты, второе число в этом случае нечётное.

Предположим сначала, что \texttt{p} нечётно, а \texttt{q} чётно. Разрешённый шаг

\begin{verbatim}
(n,k) = d(p,q)
\end{verbatim}

означает переход из \texttt{(0,0)} в \texttt{(n,k)=(dp,dq)}. Равенство
\texttt{(n,k)=d(p,q)} здесь является покоординатным равенством векторов:
\texttt{n=dp} и \texttt{k=dq}.

Так как \texttt{p} нечётно, а \texttt{q} чётно, числа \texttt{(p-1)/2} и \texttt{q/2} целые. После
шага \texttt{(n,k)} пройдём в обратную сторону построенные ранее перемещения
\texttt{(2d,0)} и \texttt{(0,2d)} указанное число раз. Получаем

\begin{verbatim}
(n,k) - ((p-1)/2)(2d,0) - (q/2)(0,2d)
= (dp,dq) - (d(p-1),0) - (0,dq)
= (d,0).
\end{verbatim}

Следовательно, \texttt{(d,0)} достижима.

Для построения \texttt{(0,d)} начнём с другого разрешённого шага:

\begin{verbatim}
(k,n) = d(q,p)
\end{verbatim}

и выполним

\begin{verbatim}
(k,n) - (q/2)(2d,0) - ((p-1)/2)(0,2d)
= (dq,dp) - (dq,0) - (0,d(p-1))
= (0,d).
\end{verbatim}

Если \texttt{p} чётно, а \texttt{q} нечётно, используются аналогичные построения:

\begin{verbatim}
(k,n) - ((q-1)/2)(2d,0) - (p/2)(0,2d) = (d,0),

(n,k) - (p/2)(2d,0) - ((q-1)/2)(0,2d) = (0,d).
\end{verbatim}

Если \texttt{d} делит \texttt{x} и \texttt{y}, существуют целые \texttt{u=x/d} и \texttt{v=y/d}. Теперь
требуемая точка строится как сумма уже построенных перемещений:

\begin{verbatim}
(x,y) = u(d,0) + v(0,d).
\end{verbatim}

Отрицательный коэффициент означает прохождение соответствующего пути в
обратную сторону. Следовательно, в этом случае делимости обеих координат
на \texttt{d} достаточно.

\subsubsection{\texorpdfstring{Случай 2: \texttt{p} и \texttt{q} оба нечётные}{Случай 2: p и q оба нечётные}}

Пусть

\begin{verbatim}
x = du,
y = dv,
\end{verbatim}

где \texttt{u} и \texttt{v} имеют одинаковую чётность.

\begin{quote}
\textbf{Именно здесь используется дополнительное условие теоремы:} при
нечётных \texttt{p} и \texttt{q} числа \texttt{u\ =\ x/d} и \texttt{v\ =\ y/d} должны иметь одинаковую
чётность. Благодаря этому остаются только два следующих случая.
\end{quote}

Если \texttt{u} и \texttt{v} чётные, числа \texttt{u/2} и \texttt{v/2} целые. Поэтому

\begin{verbatim}
(x,y) = (u/2)(2d,0) + (v/2)(0,2d).
\end{verbatim}

Правая часть является суммой уже построенных перемещений.

Если \texttt{u} и \texttt{v} нечётные, сначала сделаем один шаг

\begin{verbatim}
(n,k) = d(p,q).
\end{verbatim}

То есть после одного шага мы окажемся в \texttt{(n,k)=(dp,dq)}. Числа \texttt{u-p} и
\texttt{v-q} чётные, поскольку \texttt{u,v,p,q} нечётные. Поэтому числа \texttt{(u-p)/2} и
\texttt{(v-q)/2} целые. Добавим соответствующие целые кратные уже построенных
перемещений:

\begin{verbatim}
(n,k) + ((u-p)/2)(2d,0) + ((v-q)/2)(0,2d)
= (dp,dq) + (d(u-p),0) + (0,d(v-q))
= (du,dv)
= (x,y).
\end{verbatim}

Если один из коэффициентов отрицательный, соответствующий путь проходится
в обратную сторону.

Таким образом, указанные условия всегда достаточны. Теорема доказана.

\section{Классы достижимости}

Пусть

\begin{verbatim}
d = gcd(n,k),
p = n/d,
q = k/d,
\end{verbatim}

и числа \texttt{p} и \texttt{q} нечётные.

Точки \texttt{(x1,y1)} и \texttt{(x2,y2)} находятся в одном классе достижимости тогда
и только тогда, когда

\begin{verbatim}
∃ i | x1 ≡ x2 + i·d (mod 2d) && y1 ≡ y2 + i·d (mod 2d).
\end{verbatim}

\section{Общий ребёнок двух вершин находится на следующем уровне}

Рассматривается граф с корнем \texttt{(0,0,0,0)}. Вершина записывается как
\texttt{P\ =\ (a,b,A,B)} и представляет точку плоскости
\texttt{(an\ +\ bk,\ An\ +\ Bk)}.

Уровень вершины --- длина кратчайшего пути от корня до неё:

\begin{verbatim}
level(a,b,A,B) = max(|a|,|B|) + max(|b|,|A|).
\end{verbatim}

\subsection{Почему формула уровня верна}

Каждый шаг \texttt{D(n,k)} имеет один из двух видов:

\begin{verbatim}
(a,b,A,B) -> (a ± 1,b,A,B ± 1),
\end{verbatim}

или

\begin{verbatim}
(a,b,A,B) -> (a,b ± 1,A ± 1,B).
\end{verbatim}

Сначала рассмотрим шаги первого вида. Они изменяют только пару \texttt{(a,B)}.
За один такой шаг каждая из координат \texttt{a} и \texttt{B} может измениться не более
чем на 1. Поэтому для достижения \texttt{(a,B)} требуется не меньше

\begin{verbatim}
max(|a|,|B|)
\end{verbatim}

шагов первого вида.

Покажем, что для достижимой вершины этого числа шагов достаточно. Обозначим

\begin{verbatim}
m = max(|a|,|B|).
\end{verbatim}

У достижимой вершины \texttt{a} и \texttt{B} имеют одинаковую чётность. Число \texttt{m}
совпадает по чётности и с \texttt{a}, и с \texttt{B}. Кроме того, \texttt{\textbar{}a\textbar{}\ \textless{}=\ m} и
\texttt{\textbar{}B\textbar{}\ \textless{}=\ m}.

Любое целое число \texttt{x}, для которого \texttt{\textbar{}x\textbar{}\ \textless{}=\ m} и чётность \texttt{x} совпадает
с чётностью \texttt{m}, можно представить как сумму \texttt{m} чисел, каждое из которых
равно \texttt{+1} или \texttt{-1}. Поэтому можно выбрать знаки в \texttt{m} шагах так, чтобы
изменения первой координаты в сумме дали \texttt{a}, а изменения второй
координаты независимо дали \texttt{B}.

Следовательно, минимальное количество шагов первого вида равно

\begin{verbatim}
max(|a|,|B|).
\end{verbatim}

Точно так же шаги второго вида изменяют только пару \texttt{(b,A)}. Для
достижимой вершины \texttt{b} и \texttt{A} имеют одинаковую чётность, поэтому
минимальное количество шагов второго вида равно

\begin{verbatim}
max(|b|,|A|).
\end{verbatim}

Шаг первого вида не изменяет \texttt{(b,A)}, а шаг второго вида не изменяет
\texttt{(a,B)}. Поэтому необходимые количества шагов складываются. Получаем

\begin{verbatim}
level(a,b,A,B) = max(|a|,|B|) + max(|b|,|A|).
\end{verbatim}

Вершина \texttt{C} называется ребёнком вершины \texttt{P}, если:

\begin{enumerate}
\def\labelenumi{\arabic{enumi}.}
\tightlist
\item
  \texttt{C} и \texttt{P} являются соседями;
\item
  \texttt{level(C)\ \textgreater{}\ level(P)}.
\end{enumerate}

Здесь заранее не предполагается, что разница уровней равна 1.

\subsection{Лемма}

Если две достижимые вершины являются соседями, то их уровни отличаются
ровно на 1.

\subsubsection{Доказательство}

Пусть \texttt{P} и \texttt{C} --- соседи.

Если к кратчайшему пути от корня до \texttt{P} добавить ребро из \texttt{P} в \texttt{C},
получится путь до \texttt{C} длины \texttt{level(P)\ +\ 1}. Поэтому

\begin{verbatim}
level(C) <= level(P) + 1.
\end{verbatim}

Аналогично, если к кратчайшему пути до \texttt{C} добавить ребро из \texttt{C} в \texttt{P},
то

\begin{verbatim}
level(P) <= level(C) + 1.
\end{verbatim}

Следовательно,

\begin{verbatim}
|level(C) - level(P)| <= 1.
\end{verbatim}

Остаётся исключить равенство уровней.

Каждый шаг \texttt{D(n,k)} имеет один из двух видов:

\begin{verbatim}
(a,b,A,B) -> (a ± 1,b,A,B ± 1),
\end{verbatim}

или

\begin{verbatim}
(a,b,A,B) -> (a,b ± 1,A ± 1,B).
\end{verbatim}

В первом случае меняется чётность \texttt{a}, а во втором --- чётность \texttt{b}.
Поэтому каждый шаг меняет чётность числа \texttt{a\ +\ b}.

В корне \texttt{a\ +\ b\ =\ 0}. Следовательно, для каждой достижимой вершины
чётность её уровня совпадает с чётностью \texttt{a\ +\ b}. У соседних вершин
чётность \texttt{a\ +\ b} различна, поэтому их уровни не могут быть равны.

Вместе с ранее полученным неравенством это даёт

\begin{verbatim}
|level(C) - level(P)| = 1.
\end{verbatim}

Лемма доказана.

\subsection{Утверждение}

Если две вершины \texttt{P1} и \texttt{P2} имеют общего ребёнка \texttt{C}, то \texttt{P1} и \texttt{P2}
находятся на одном уровне.

\subsubsection{Доказательство}

Так как \texttt{C} --- ребёнок \texttt{P1}, вершины \texttt{C} и \texttt{P1} являются соседями и
\texttt{level(C)\ \textgreater{}\ level(P1)}. По доказанной лемме их уровни отличаются ровно
на 1. Значит,

\begin{verbatim}
level(C) = level(P1) + 1.
\end{verbatim}

Аналогично, поскольку \texttt{C} --- ребёнок \texttt{P2},

\begin{verbatim}
level(C) = level(P2) + 1.
\end{verbatim}

Следовательно,

\begin{verbatim}
level(P1) + 1 = level(P2) + 1,
\end{verbatim}

откуда

\begin{verbatim}
level(P1) = level(P2).
\end{verbatim}

Утверждение доказано.

\section{У каждой ненулевой вершины не более одного одинокого ребёнка}

Рассматривается граф с корнем \texttt{(0,0,0,0)}. Для достижимой вершины

\begin{verbatim}
P = (a,b,A,B)
\end{verbatim}

уровень равен

\begin{verbatim}
level(P) = max(|a|,|B|) + max(|b|,|A|).
\end{verbatim}

Соседний шаг имеет один из двух видов:

\begin{verbatim}
(a,b,A,B) -> (a ± 1,b,A,B ± 1),
\end{verbatim}

или

\begin{verbatim}
(a,b,A,B) -> (a,b ± 1,A ± 1,B).
\end{verbatim}

Родителем вершины \texttt{P} называется её сосед на меньшем уровне. Одиноким
ребёнком называется вершина, имеющая ровно одного родителя.

Будем использовать следующие обозначения. Для ненулевого целого \texttt{t}
через \texttt{t⁻} обозначим движение на 1 к нулю:

\begin{verbatim}
t⁻ = t-1, если t>0,
t⁻ = t+1, если t<0.
\end{verbatim}

Через \texttt{t⁺} обозначим движение на 1 от нуля:

\begin{verbatim}
t⁺ = t+1, если t>0,
t⁺ = t-1, если t<0.
\end{verbatim}

Для \texttt{t=0} существуют два движения от нуля: к \texttt{1} и к \texttt{-1}.

\subsection{Лемма: какие вершины имеют ровно одного родителя}

Достижимая ненулевая вершина имеет ровно одного родителя тогда и только
тогда, когда она имеет один из видов

\begin{verbatim}
(a,0,0,B), где |a|=|B|>0,
\end{verbatim}

или

\begin{verbatim}
(0,b,A,0), где |b|=|A|>0.
\end{verbatim}

\subsubsection{Доказательство}

Разобьём доказательство на два случая.

\subsubsection{\texorpdfstring{Случай 1: одна из пар \texttt{(a,B)} и \texttt{(b,A)} равна нулю}{Случай 1: одна из пар (a,B) и (b,A) равна нулю}}

Без ограничения общности предположим

\begin{verbatim}
b=A=0,
\end{verbatim}

а пара \texttt{(a,B)} ненулевая.

\paragraph{\texorpdfstring{Случай 1.1: \texttt{\textbar{}a\textbar{}=\textbar{}B\textbar{}}}{Случай 1.1: \textbar a\textbar=\textbar B\textbar{}}}

Обозначим общее значение через \texttt{m\textgreater{}0}. Тогда вклад пары \texttt{(a,B)} в уровень
равен \texttt{m}.

Чтобы уменьшить уровень, необходимо уменьшить и \texttt{\textbar{}a\textbar{}}, и \texttt{\textbar{}B\textbar{}}.
Существует ровно один такой сосед:

\begin{verbatim}
(a⁻,0,0,B⁻).
\end{verbatim}

Его уровень на 1 меньше:

\begin{verbatim}
max(|a⁻|,|B⁻|) = m-1.
\end{verbatim}

В каждом из остальных трёх соседей хотя бы одна из величин \texttt{\textbar{}a\textbar{}}, \texttt{\textbar{}B\textbar{}}
увеличивается до \texttt{m+1}. Поэтому эти соседи находятся на большем уровне.
Следовательно, вершина имеет ровно одного родителя.

\paragraph{\texorpdfstring{Случай 1.2: \texttt{\textbar{}a\textbar{}\ !=\ \textbar{}B\textbar{}}}{Случай 1.2: \textbar a\textbar{} != \textbar B\textbar{}}}

Без ограничения общности предположим

\begin{verbatim}
|a|>|B|.
\end{verbatim}

Так как вершина достижима, \texttt{a} и \texttt{B} имеют одинаковую чётность. Значит,
\texttt{\textbar{}a\textbar{}} и \texttt{\textbar{}B\textbar{}} также имеют одинаковую чётность. Их положительная разность
не может равняться 1, поэтому

\begin{verbatim}
|a| >= |B|+2.
\end{verbatim}

Уменьшим \texttt{\textbar{}a\textbar{}} на 1. Координату \texttt{B} можно изменить обоими разрешёнными
способами. Если \texttt{B\ !=\ 0}, получаются соседи

\begin{verbatim}
(a⁻,0,0,B⁻)
\end{verbatim}

и

\begin{verbatim}
(a⁻,0,0,B⁺).
\end{verbatim}

Для обоих выполняется

\begin{verbatim}
|B⁻| <= |a|-1,
|B⁺| = |B|+1 <= |a|-1.
\end{verbatim}

Поэтому уровень обоих соседей равен \texttt{\textbar{}a\textbar{}-1}, то есть оба являются
родителями.

Если \texttt{B=0}, вместо \texttt{B⁻} и \texttt{B⁺} берутся два возможных значения \texttt{-1} и
\texttt{1}. Неравенство \texttt{\textbar{}a\textbar{}\textgreater{}=2} снова показывает, что оба полученных соседа ---
родители.

Следовательно, при \texttt{\textbar{}a\textbar{}!=\textbar{}B\textbar{}} родителей не меньше двух.

Случай, в котором \texttt{(a,B)=(0,0)}, а \texttt{(b,A)} ненулевая, доказывается точно
так же. В нём единственный родитель существует ровно при \texttt{\textbar{}b\textbar{}=\textbar{}A\textbar{}\textgreater{}0}.

\subsubsection{\texorpdfstring{Случай 2: обе пары \texttt{(a,B)} и \texttt{(b,A)} ненулевые}{Случай 2: обе пары (a,B) и (b,A) ненулевые}}

Покажем, что каждая ненулевая пара даёт хотя бы одного родителя.

Для пары \texttt{(a,B)} возможны три варианта:

\begin{enumerate}
\def\labelenumi{\arabic{enumi}.}
\tightlist
\item
  Если \texttt{\textbar{}a\textbar{}=\textbar{}B\textbar{}\textgreater{}0}, движение \texttt{(a,B)-\textgreater{}(a⁻,B⁻)} уменьшает первый максимум
  на 1.
\item
  Если \texttt{\textbar{}a\textbar{}\textgreater{}\textbar{}B\textbar{}}, то из достижимости следует \texttt{\textbar{}a\textbar{}\textgreater{}=\textbar{}B\textbar{}+2}. Движение
  координаты \texttt{a} к нулю и любое допустимое движение координаты \texttt{B}
  уменьшает первый максимум на 1.
\item
  Если \texttt{\textbar{}B\textbar{}\textgreater{}\textbar{}a\textbar{}}, применяется симметричное движение: \texttt{B} движется к
  нулю, а \texttt{a} выбирается так, чтобы первый максимум уменьшился на 1.
\end{enumerate}

Так получается родитель, отличающийся от \texttt{P} только в координатах
\texttt{a} и \texttt{B}.

Такие же три варианта для ненулевой пары \texttt{(b,A)} дают родителя,
отличающегося от \texttt{P} только в координатах \texttt{b} и \texttt{A}.

Эти два родителя различны: первый изменяет \texttt{(a,B)}, а второй --- \texttt{(b,A)}.
Следовательно, если обе пары ненулевые, у вершины не меньше двух
родителей.

Лемма доказана.

\subsection{Утверждение}

У каждой вершины, отличной от \texttt{(0,0,0,0)}, не более одного одинокого
ребёнка.

\subsubsection{Доказательство}

По лемме все одинокие дети имеют один из видов

\begin{verbatim}
(εr,0,0,δr)
\end{verbatim}

или

\begin{verbatim}
(0,εr,δr,0),
\end{verbatim}

где \texttt{r\textgreater{}=1}, а \texttt{ε,δ} независимо принимают значения \texttt{-1} и \texttt{1}.

Если \texttt{r\textgreater{}=2}, единственный родитель первой вершины равен

\begin{verbatim}
(ε(r-1),0,0,δ(r-1)).
\end{verbatim}

По этому родителю однозначно восстанавливаются \texttt{r}, \texttt{ε} и \texttt{δ}.
Следовательно, ненулевая вершина не может быть родителем двух различных
одиноких детей первого вида.

Для второго вида единственный родитель равен

\begin{verbatim}
(0,ε(r-1),δ(r-1),0),
\end{verbatim}

и применяется то же рассуждение.

Родитель одинокого ребёнка первого вида может совпасть с родителем
одинокого ребёнка второго вида только тогда, когда все четыре его
координаты равны нулю. Это происходит при \texttt{r=1}.

Таким образом, только корень \texttt{(0,0,0,0)} является общим родителем
нескольких одиноких детей. Любая другая вершина имеет не более одного
одинокого ребёнка. Утверждение доказано.

\section{План доказательства классификации бинарных отношений}

Для \texttt{k\ ≥\ n\ ≥\ 0} обозначим через \texttt{D\_(k,n)(x,y)} отношение

\begin{verbatim}
y - x ∈ {(±k,±n),(±n,±k)}.
\end{verbatim}

Рассматривается определимость без параметров: в формуле не используются
заранее выбранные точки решётки.

Предполагаемое утверждение: каждое бинарное отношение, определимое в
\texttt{\textless{}Z²,E\textgreater{}}, является либо конечным объединением отношений \texttt{D\_(k,n)}, либо
дополнением конечного объединения таких отношений.

\subsection{План доказательства}

\begin{enumerate}
\def\labelenumi{\arabic{enumi}.}
\item
  Для каждого фиксированного \texttt{k,n} отношение \texttt{D\_(k,n)} определимо через \texttt{E}.
  Это даёт определимость любого конечного объединения таких отношений и его
  дополнения.
\item
  Переносы, повороты на угол, кратный \texttt{90°}, и отражения сохраняют \texttt{E}.
  Поэтому они сохраняют любое отношение, определимое через \texttt{E}. Любые две пары
  точек одного типа \texttt{D\_(k,n)} переводятся друг в друга одним из таких
  преобразований. Следовательно, определимая формула имеет постоянное значение
  на каждом \texttt{D\_(k,n)}.
\item
  Пусть бинарное отношение задаётся FO-формулой \texttt{φ(x,y)}. По
  \href{https://homepages.inf.ed.ac.uk/libkin/fmt/fmt.pdf}{теореме Гайфмана 4.22 в книге Леонида Либкина \emph{Elements of Finite Model Theory}}
  эта формула равносильна конечной булевой комбинации:

  \begin{itemize}
  \tightlist
  \item
    локальных формул вокруг \texttt{(x,y)};
  \item
    предложений без свободных переменных, утверждающих существование нескольких
    достаточно удалённых точек с заданным локальным свойством.
  \end{itemize}
\item
  Для фиксированной решётки предложения без свободных переменных имеют
  постоянные значения и не различают пары \texttt{(x,y)}. Поскольку булева комбинация
  конечна, среди радиусов всех её локальных формул существует максимальный
  радиус \texttt{r}.
\item
  Если \texttt{d(x,y)\ \textgreater{}\ 2r+1}, шары радиуса \texttt{r} вокруг \texttt{x} и \texttt{y} не пересекаются и
  между ними нет рёбер. Поэтому окрестность пары \texttt{(x,y)} является объединением
  двух отдельных одинаковых шаров с отмеченными центрами.
\item
  Окрестности радиуса \texttt{r} любых двух пар, расстояние внутри которых больше
  \texttt{2r+1}, изоморфны с сохранением отмеченных \texttt{x} и \texttt{y}. Все локальные части
  формулы, а значит и \texttt{φ}, имеют на таких парах одно общее значение.
\item
  Для пары типа \texttt{D\_(k,n)} расстояние по клеткам равно \texttt{k+n}. При
  \texttt{k+n\ ≤\ 2r+1} существует лишь конечное число пар целых чисел
  \texttt{k\ ≥\ n\ ≥\ 0}. Поэтому отдельно формула может выбирать только конечное число
  типов \texttt{D\_(k,n)}; на всех остальных типах её значение одинаково.
\item
  Если общее значение на достаточно далёких парах ложно, определяемое
  отношение является конечным объединением \texttt{D\_(k,n)}. Если оно истинно,
  дополнение определяемого отношения является конечным объединением
  \texttt{D\_(k,n)}.
\end{enumerate}

\subsection{Отношения большей арности}

Предыдущее утверждение непосредственно классифицирует только бинарные
отношения. Для отношения

\begin{verbatim}
R(x₁,...,xₘ)
\end{verbatim}

нужно учитывать не только расстояние между одной парой точек, но и взаимное
расположение всех точек кортежа.

Полезный образец даёт работа А. Л. Семёнова и С. Ф. Сопрунова
\href{https://arxiv.org/abs/1201.4439}{\emph{Lattice of relational algebras definable in integers with successor}}
(\href{https://www.mathnet.ru/eng/im9107}{опубликованная версия}). В её лемме 6
доказывается аналог следующего принципа: у каждого определимого отношения есть
конечная ширина. После некоторой границы формула перестаёт различать точные
расстояния и видит только то, какие элементы расположены близко друг к другу.

Приведённый ниже план является переносом этой идеи на \texttt{\textless{}Z²,E\textgreater{}}, а не
утверждением из самой статьи.

\subsubsection{1. Лемма о ширине}

Для каждого определимого отношения \texttt{R(x₁,...,xₘ)} нужно доказать существование
числа \texttt{w} со следующим свойством. Если отмеченные \texttt{w}-окрестности двух кортежей

\begin{verbatim}
(x₁,...,xₘ) и (y₁,...,yₘ)
\end{verbatim}

изоморфны, причём каждая отмеченная точка \texttt{xᵢ} переходит в \texttt{yᵢ}, то

\begin{verbatim}
R(x₁,...,xₘ) ⇔ R(y₁,...,yₘ).
\end{verbatim}

Для структуры \texttt{\textless{}Z²,E\textgreater{}} это следует из теоремы Гайфмана: предложения без
свободных переменных имеют в фиксированной решётке постоянное значение, а среди
радиусов локальных частей формулы можно выбрать наибольший.

\subsubsection{2. Разбиение переменных на близкие группы}

Для кортежа \texttt{(x₁,...,xₘ)} построим конечный граф на номерах переменных:

\begin{verbatim}
i соединено с j ⇔ d(xᵢ,xⱼ) ≤ 2w+1.
\end{verbatim}

Компоненты связности этого графа назовём близкими группами. Точки из разных
групп имеют непересекающиеся \texttt{w}-окрестности, между которыми нет рёбер.

Например, для трёх переменных возможны разбиения

\begin{verbatim}
{x}, {y}, {z};
{x,y}, {z};
{x,z}, {y};
{y,z}, {x};
{x,y,z}.
\end{verbatim}

\subsubsection{3. Тип конфигурации целой группы}

Для двух точек тип задаётся одним отношением \texttt{D\_(k,n)}. Для трёх и более точек
нужно описывать всю конфигурацию одновременно. Нельзя независимо выбирать
поворот или отражение для каждой пары: одно геометрическое преобразование должно
действовать сразу на всю близкую группу.

Пусть \texttt{v₂,...,vₛ} --- фиксированные векторы. Обозначим через
\texttt{T\_(v₂,...,vₛ)(x₁,...,xₛ)} следующий тип конфигурации:

\begin{verbatim}
существует один поворот или отражение g такой, что
xᵢ - x₁ = g(vᵢ) для всех i = 2,...,s.
\end{verbatim}

При \texttt{s=2} получается бинарный тип \texttt{D\_(k,n)}. Таким образом, отношения \texttt{T}
являются естественными многоместными аналогами \texttt{D\_(k,n)}.

Это пока геометрическое обозначение типа. Его определимость через \texttt{E} нужно
доказывать отдельно для каждой фиксированной конечной конфигурации, описывая её
ограниченную отмеченную окрестность.

\subsubsection{4. Почему типов конечное число}

В одной близкой группе находится не более \texttt{m} точек. Между соседними точками
графа группы расстояние не превосходит \texttt{2w+1}. Поэтому диаметр группы не
превосходит

\begin{verbatim}
(m-1)(2w+1).
\end{verbatim}

В шаре ограниченного радиуса решётки находится конечное число точек.
Следовательно, для фиксированных \texttt{m} и \texttt{w} существует только конечное число
типов близких отмеченных конфигураций с точностью до переноса, поворота и
отражения.

Отсюда ожидается следующая нормальная форма для любого \texttt{m}-местного отношения:

\begin{verbatim}
разбиение переменных на близкие группы
+ тип конфигурации каждой группы
+ условие, что разные группы достаточно далеки.
\end{verbatim}

Само отношение является конечным объединением таких случаев.

\subsubsection{5. Переход к решётке определимости}

Нормальная форма описывает отдельное отношение, но ещё не определяет его место
в решётке определимости. Следующий этап по образцу Семёнова и Сопрунова:

\begin{enumerate}
\def\labelenumi{\arabic{enumi}.}
\tightlist
\item
  Зафиксировать произвольное определимое отношение \texttt{R}.
\item
  Выбрать для него ширину \texttt{w}.
\item
  Определить, какие конечные конфигурации различает \texttt{R}.
\item
  Рассмотреть перестановки точек, сохраняющие \texttt{R}.
\item
  Проверить, сохраняют ли эти перестановки некоторое бинарное отношение
  \texttt{D\_(k,n)}.
\item
  Если подходящего \texttt{D\_(k,n)} нет, искать новый действительно многоместный
  базовый тип отношения.
\end{enumerate}

Поэтому пока нельзя утверждать, что каждое отношение большей арности взаимно
определимо с некоторым бинарным \texttt{D\_(k,n)}. Это отдельное утверждение, которое
нужно доказать или опровергнуть.


\begingroup
\renewcommand{\refname}{Список использованных источников}
\let\section\subsection
\addcontentsline{toc}{subsection}{\refname}
\begin{thebibliography}{9}

\bibitem{semenov-soprunov}
Семёнов, А.~Л. Решётка определимости (редуктов) для целых чисел с операцией
следования / А.~Л. Семёнов, С.~Ф. Сопрунов. --- Текст : электронный~//
Известия Российской академии наук. Серия математическая. --- 2021. ---
Т.~85, №~6. --- С.~245--258. --- DOI~10.4213/im9107. ---
URL: \url{https://www.mathnet.ru/rus/im9107}
(дата обращения: 19.08.2026).

\bibitem{libkin-fmt}
Libkin, L. Elements of Finite Model Theory / L.~Libkin. ---
Berlin ; Heidelberg : Springer, 2004. --- Текст : электронный. ---
DOI~10.1007/978-3-662-07003-1. ---
URL: \url{https://homepages.inf.ed.ac.uk/libkin/fmt/fmt.pdf}
(дата обращения: 19.08.2026).

\end{thebibliography}
\endgroup

\end{document}
>>>>>>> 8ed2429 (Moved things to latex)
